\documentclass[11pt,a4paper]{article}
\usepackage[utf8]{inputenc}
\usepackage[T1]{fontenc}
\usepackage[margin=2.25cm]{geometry}
\usepackage{amsmath,amssymb}
\usepackage{array}
\usepackage{booktabs}
\usepackage{enumitem}
\usepackage{xcolor}
\usepackage[numbers]{natbib}
\usepackage[hidelinks]{hyperref}
\usepackage{parskip}
\usepackage{xurl}
\usepackage{microtype}
\usepackage{listings}
\usepackage{multicol}

\sloppy
\definecolor{bettergrey}{gray}{0.95}
\definecolor{betterblue}{RGB}{33, 65, 196}
\newcommand{\field}[1]{\texttt{#1}}
\newcommand{\stage}[2]{\subsection{\textcolor{betterblue}{#1}: #2}}

\lstdefinestyle{promptstyle}{
  basicstyle=\ttfamily\small,
  breaklines=true,
  frame=single,
  backgroundcolor=\color{bettergrey},
  columns=fullflexible,
  keepspaces=true
}

\title{\textbf{Synthetic Dataset Construction Pipeline}}
\date{}
\author{}
\begin{document}
\maketitle
\tableofcontents
% \newpage

\section{Pipeline Core Design}

The new pipeline reverses the old MIMIC-IV workflow. The new benchmark starts from a structured plan:

\begin{enumerate}[noitemsep, topsep=0pt]
	\item define query facets and answer structure
	\item generate atomic clinical facts that support each facet
	\item render those facts into narrative clinical text chunks
	\item generate redundant gold clusters and hard distractor clusters
	\item generate natural multi-aspect queries from the same hidden plan
	\item derive qrels and gold answers directly from the fact--chunk--facet mapping
	\item validate the embedding geometry before final evaluation.
\end{enumerate}

\subsection{Practical MVP}

Do not start with every idea in the full design. The first implementation should be small and deterministic:

\begin{itemize}[noitemsep, topsep=0pt]
	\item 4 conditions;
	\item 2 clinical axes: treatment duration and rehabilitation outcome;
	\item 4 facets per query;
	\item 10 gold chunks per facet;
	\item 30 hard distractors per query;
	\item 100--150 queries;
	\item top-$k$, MMR, and facility-location only.
\end{itemize}

No Prolog, no agentic retrieval, no information-not-found questions, no conflicting-note resolution in the first version. Treat \field{logical\_form} as JSON stored in \field{query\_plans.parquet}. The ``solver'' is just pandas filters over structured columns.

\begingroup
\renewcommand{\arraystretch}{1.15}
\begin{center}
	\begin{tabular}{@{}p{3.0cm}p{11.2cm}@{}}
		\toprule
		\textbf{File/script} & \textbf{What to implement first} \\
		\midrule
		\field{ontology.yaml} & Hand-written conditions, subgroups, allowed treatments, duration bins, rehab-outcome bins. \\
		\field{01\_plans.py} & Nested loops over condition and subgroup pairs. Create 4 facets: subgroup A duration, subgroup A outcome, subgroup B duration, subgroup B outcome. \\
		\field{02\_facts.py} & Expand each facet into structured rows. One fact row supports one facet only. \\
		\field{03\_chunks.py} & Render facts with deterministic templates first. Add LLM rewriting only after metrics work. \\
		\field{04\_queries.py} & Use one query template and one answer template. LLM paraphrasing is optional and must not change labels. \\
		\field{05\_qrels.py} & Join chunks to facts. Gold rows get relevance 1; distractors get relevance 0. \\
		\field{06\_embed.py} & Embed chunks and queries. Store vectors in LanceDB or parquet. \\
		\field{07\_geometry.py} & Keep only queries where all facets appear in top-$N$, top-$k$ over-samples one facet, and distractors appear in the pool. \\
		\field{08\_evaluate.py} & Compute Facet Coverage, Gold Precision, Distractor Rate, and Dominant-Cluster Concentration. \\
		\bottomrule
	\end{tabular}
\end{center}
\endgroup

Concrete keep rules for the MVP:

\begin{itemize}[noitemsep, topsep=0pt]
	\item candidate pool: $N=300$ query-local or $N=1000$ full-corpus;
	\item every required facet has at least one gold chunk in top-$N$;
	\item top-10 baseline retrieves at least 5 chunks from one facet;
	\item mean in-facet similarity exceeds mean cross-facet similarity by at least 0.03;
	\item at least 10 distractors appear in top-$N$.
\end{itemize}

% \subsection{Running Example}

% \textbf{Strategy.} Generate a structured plan before generating chunks or queries. A clinical query should be backed by a graph or tree of candidate facets, such as primary condition, demographic subgroup, comorbidity subgroup, treatment duration, complications, rehabilitation outcome, discharge disposition, and follow-up needs.

% For example, an encephalitis/myelitis query family could start from the following candidate facet pool:

% \begin{itemize}[noitemsep, topsep=0pt]
% 	\item primary condition: encephalitis or myelitis
% 	\item demographic candidates: age $>75$, age 40--64
% 	\item comorbidity candidates: uncomplicated diabetes, renal disease, immunosuppression
% 	\item treatment candidates: antiviral duration, steroid exposure, delayed treatment start
% 	\item outcome candidates: rehabilitation outcome, discharge disposition, persistent neurologic deficit, follow-up needs
% 	\item distractor candidates: same condition but wrong subgroup, same subgroup but wrong condition, same treatment vocabulary but wrong disease.
% \end{itemize}

% The seven steps above then instantiate as follows.

% \begin{enumerate}[leftmargin=*, itemsep=5pt, topsep=0pt]
% 	\item \textbf{Query facets and answer structure.} The hidden plan selects four compatible candidates and turns them into required query facets:
% 	      \field{F1}: diabetes + treatment duration
% 	      \field{F2}: diabetes + rehabilitation outcome
% 	      \field{F3}: age $>75$ + treatment duration
% 	      \field{F4}: age $>75$ + rehabilitation outcome.
% 	      The expected answer must compare the diabetes subgroup against the elderly subgroup across both clinical axes.

% 	\item \textbf{Atomic clinical facts.} The generator creates structured facts before writing any text. For example:
% 	      \begin{lstlisting}[style=promptstyle]
% {
%   "fact_id": "FACT_001",
%   "primary_condition": "encephalitis",
%   "age_group": "age_40_64",
%   "comorbidities": ["uncomplicated_diabetes"],
%   "intervention": "intravenous acyclovir",
%   "treatment_duration": "10 days",
%   "rehab_disposition": "home physical therapy",
%   "outcome": "near-baseline mobility at discharge",
%   "facet_ids": ["F1", "F2"],
%   "must_not_mention": ["age over 75", "inpatient rehabilitation"]
% }
% 	      \end{lstlisting}

% 	\item \textbf{Narrative clinical chunks.} The fact is rendered into one or more retrieval chunks while preserving the labels:
% 	      \begin{lstlisting}[style=promptstyle]
% C_001: The patient was treated for encephalitis with a 10-day course of
% intravenous acyclovir. History included uncomplicated diabetes without
% documented end-organ complications. By discharge the patient had returned
% close to baseline mobility and was referred for home physical therapy.

% metadata: fact_ids=["FACT_001"], facet_ids=["F1","F2"], cluster_id="C_diab_rehab"
% 	      \end{lstlisting}

% 	\item \textbf{Gold and distractor clusters.} The benchmark generates several chunks per facet, for example 10 chunks for each of \field{F1}--\field{F4}. It also generates hard distractors such as elderly stroke rehabilitation notes, encephalitis patients without diabetes, or diabetes patients with pneumonia treatment durations.

% 	\item \textbf{Natural query.}
% 	      \begin{quote}
% 		      For patients diagnosed with encephalitis or myelitis, how do treatment durations and rehabilitative outcomes differ between patients with uncomplicated diabetes and patients older than 75?
% 	      \end{quote}

% 	\item \textbf{Qrels and gold answer.} The qrels are derived directly from the hidden mapping, e.g. \field{(Q\_001, C\_001, F1, relevant)} and \field{(Q\_001, C\_001, F2, relevant)}. The gold answer is then synthesized from all required facets, with citations to one or more chunks per facet.

% 	\item \textbf{Geometry validation.} After embedding, the query is kept only if all four facets appear in the top-$N$ candidate pool, top-$k$ over-samples at least one dominant cluster, and the distractors are close enough to challenge retrieval without becoming gold evidence.
% \end{enumerate}

\subsection{Literature overview}

\subsubsection{Generic papers}
The papers start from an already existing data source and build a QA benchmark on top of it. The main ideas in the papers are: breaking down the data source into structured data supporting  knowledge discovery and query generation, allowing for control of query difficulty, explicit validation of the produced synthetic data, explicit evidence provenance, different methods to ensure and validate the diversity of the generated data.\\
This pipeline borrows their construction methods, but changes the starting point: MIMIC-IV is used only as a reference for terminology, condition families, comorbidity co-occurrences, demographic bins, and note style, but the real starting point is the structured metadata about the dataset we want to generate. The final corpus is generated from a hidden schema so that the retrieval problem has the multi-cluster structure that the raw discharge note corpus did not provide, so it solves by construction many of the challenges encountered in the papers.

Useful ideas from the literature:

\begin{itemize}[noitemsep, topsep=0pt]
	\item \textbf{Evidence-first construction.} MHTS builds questions from claims, clusters, evidence mappings, and difficulty signals \citep{lee2025mhts}. DRAGON stores explicit links among documents, queries, clues, answers, and citations \citep{shen2026dragon}. In this pipeline, those ideas become \field{query\_plans.parquet}, \field{clinical\_facts.parquet}, \field{chunks.parquet}, and \field{qrels.parquet}. The LLM writes the clinical text, but the gold labels come from the hidden mapping.
	\item \textbf{Multi-aspect query structure.} ClaimSpect and MQ4CS are useful because they treat a broad user query as a set of separable aspects \citep{kargupta2025claimspect,abbasiantaeb2026multiaspect}. Here that decomposition is used before generation: a query plan explicitly chooses facets such as subgroup, treatment duration, rehabilitation outcome, and discharge disposition.
	\item \textbf{Question-aware corpus construction.} EnterpriseRAG-Bench generates documents around planned question structures, cross-document coherence, realistic noise, completeness clusters, conflicting near-duplicates, answer facts, and conformance checks \citep{sun2026enterpriseragbench}. Here this becomes a synthetic clinical archive with shared condition families, care-path scaffolds, note-style scaffolds, intentional hard distractors, and query-local cluster geometry.
	\item \textbf{On-demand hidden worlds.} PhantomWiki generates a fresh fictional universe from hidden symbolic facts, then derives templated questions and exhaustive answers with a solver \citep{gong2025phantomwiki}. Here the equivalent is a generated clinical world with fixed seeds, hidden fact tables, optional query templates/logical forms, and qrels derived before LLM narrative rendering.
	\item \textbf{Difficulty control.} MHTS motivates storing difficulty as number of facets, number of evidence pieces, and query--evidence similarity \citep{lee2025mhts}. In this benchmark, difficulty is controlled by \field{n\_facets}, \field{gold\_chunks\_per\_facet}, dominant-cluster size, distractor count, and whether complementary facets are less similar to the surface query.
	\item \textbf{Validation and provenance.} Rahmani et al. and RAGVAL support automatic test-collection construction but also show why validation matters \citep{rahmani2024synthetic,kenneweg2024ragval}. This pipeline therefore stores fact IDs, facet IDs, cluster IDs, etc.
	\item \textbf{Diversity.} DATAGEN uses attribute-guided generation, diversity-oriented example selection, and group checking to avoid duplicate synthetic examples \citep{huang2025datagen}. Driouich et al. use a diversity agent that clusters and samples source documents for broader RAG coverage \citep{driouich2025diverseprivate}. Those methods cannot be copied literally, because this benchmark needs both diversity and planned redundancy. The goal here is to preserve redundancy inside a dominant gold cluster, enforce separation among different gold facets, and create near-miss distractor clusters that are close enough to challenge MMR.
	\item \textbf{Downstream filtering.} Kim et al. argue that synthetic generators should be judged by the downstream behavior their data induces \citep{kim2025agorabench}. For the thesis, a generated item is only useful if the candidate pool actually exposes a coverage-sensitive retrieval problem.
	\item \textbf{Optional prompt optimization.} If prompts are unstable, prompt selection can be treated as a multi-objective optimization problem. Pastri{\'a}n et al. provide one example of this kind of automatic prompt evolution, using evolutionary multi-objective prompt learning rather than manually choosing one prompt by inspection \citep{pastrian2026evolutionary}. This is an optional optimization for improving the generator (maybe overkill in this phase).
\end{itemize}

\subsubsection{Clinical papers and MIMIC-IV role}

MIMIC-IV can be used as a reference for terminology, frequent conditions, comorbidity distributions, age groups, etc.

This can be strengthened with ideas from the medical papers for synthetic data construction:

\begin{itemize}[]
	\item simple care-path templates that say, for a given condition and subgroup, which treatments and outcomes are plausible. This follows the same general idea as HealthBranches \citep{cosentino2025healthbranches};
	\item a small controlled vocabulary of medical terms, so the generator does not randomly alternate between inconsistent names for the same condition, treatment, or outcome. This is the practical part of knowledge-infused medical QA systems such as KI-MAG \citep{zafar2024kimag};
	\item explicit patient-group rules, such as ``encephalitis + age over 75 + rehabilitation outcome'' or ``pneumonia + diabetes + treatment duration''. This borrows the idea of decomposing clinical questions into condition, subgroup, treatment, and outcome constraints \citep{ziletti2025cohorts}.
\end{itemize}

\section{Pipeline}

\stage{Stage 1}{Build a Facet Catalogue}

Build it manually first. MIMIC-IV can be used as a reference for common conditions, age groups, and comorbidities, but it should not drive the qrels. Care-path templates, controlled vocabulary, cohort-style criteria, and rubric-style facet checks can be added later if the generated data is not clinically coherent enough \citep{cosentino2025healthbranches,zafar2024kimag,ziletti2025cohorts,gong2026meddialogrubrics}.

Add a small archive scaffold before any chunks are generated: condition-family overviews, care-path templates, latent patient/profile tables, note/source styles, and style-agent instructions. This adapts EnterpriseRAG-Bench's scaffolding approach to clinical data \citep{sun2026enterpriseragbench}. Store fixed seeds and schema versions so the benchmark can be regenerated as fresh instances, following PhantomWiki's on-demand evaluation idea \citep{gong2025phantomwiki}.

We should have:

\begin{itemize}[noitemsep, topsep=0pt]
	\item 4--8 primary conditions to use in the benchmark;
	\item demographic groups, e.g. age $<40$, age 40--64, age 65--75, age $>75$;
	\item comorbidity groups, e.g. uncomplicated diabetes, renal disease, chronic pulmonary disease;
	\item clinical axes we want questions to ask about, e.g. treatment duration, intervention choice, complications, rehabilitation outcome, discharge disposition;
	\item allowed combinations, so the generator does not create completely implausible cases.
\end{itemize}

Example catalogue row:

\begin{lstlisting}[style=promptstyle]
condition: encephalitis_or_myelitis
allowed_subgroups:
  demographics: [age_over_75, age_40_64]
  comorbidities: [uncomplicated_diabetes, renal_disease]
allowed_axes:
  - treatment_duration
  - rehabilitation_outcome
  - discharge_disposition
allowed_treatments: [acyclovir, steroids, supportive_care]
allowed_outcomes: [home_rehab, inpatient_rehab, persistent_deficit]
\end{lstlisting}

\stage{Stage 2}{Hidden Query Plans}

Generate query plans from the catalogue before writing any query text. A query plan is the hidden recipe saying which facets must be covered. For example, choose:

\begin{itemize}[noitemsep, topsep=0pt]
	\item one primary condition: encephalitis/myelitis;
	\item two subgroups to compare: uncomplicated diabetes vs. age $>75$;
	\item two clinical axes: treatment duration and rehabilitation outcome.
\end{itemize}

This produces four required facets:

\begin{lstlisting}[style=promptstyle]
query_id: Q_001
primary_condition: encephalitis_or_myelitis
query_type: subgroup_comparison
template_id: subgroup_comparison_two_axes
logical_form:
  condition: encephalitis_or_myelitis
  arms: [uncomplicated_diabetes, age_over_75]
  axes: [treatment_duration, rehabilitation_outcome]
  operator: compare
required_facets:
  F1: uncomplicated_diabetes + treatment_duration
  F2: uncomplicated_diabetes + rehabilitation_outcome
  F3: age_over_75 + treatment_duration
  F4: age_over_75 + rehabilitation_outcome
answer_outline:
  - summarize treatment duration for uncomplicated diabetes
  - summarize rehabilitation outcome for uncomplicated diabetes
  - summarize treatment duration for age_over_75
  - summarize rehabilitation outcome for age_over_75
  - compare the two subgroups
difficulty:
  n_facets: 4
  dominant_gold_cluster: C1
  dominant_gold_chunks: 18
  complementary_gold_chunks_per_cluster: 8
  hard_distractor_chunks: 24
  expected_topk_failure: oversample_dominant_gold_cluster
split: validation
\end{lstlisting}

The output of this stage is \field{query\_plans.parquet}. This follows the answer-first idea used in synthetic RAG benchmark construction \citep{lee2025mhts,shen2026dragon}, the question-aware corpus construction of EnterpriseRAG-Bench \citep{sun2026enterpriseragbench}, and PhantomWiki's separation between hidden logical form and surface question text \citep{gong2025phantomwiki}. The important point is that diversity is already specified here: the plan says which clusters are gold, which one is intentionally redundant, which facets are complementary, and which distractor families should be generated.

\stage{Stage 3}{Generate Structured Fact Rows}

For each required facet, generate several structured fact rows. A fact row is the hidden source of truth for one synthetic case. It says what the generated text is allowed to mention and which facet it supports.

Example:

\begin{lstlisting}[style=promptstyle]
fact_id: FACT_001
query_id: Q_001
facet_ids: [F3, F4]
condition: encephalitis
age_group: age_over_75
comorbidities: [hypertension]
treatment: acyclovir
treatment_duration: 14_days
rehab_outcome: inpatient_rehab_with_persistent_gait_instability
must_mention:
  - age over 75
  - 14-day antiviral course
  - inpatient rehabilitation
  - persistent gait instability
must_not_mention:
  - uncomplicated diabetes
  - complete recovery
\end{lstlisting}

For each facet, generate enough rows to create a cluster. These rows should vary across synthetic patients while preserving the same facet meaning. The output of this stage is \field{clinical\_facts.parquet}.

For the MVP, this rule-check layer is just pandas filtering over \field{clinical\_facts.parquet}. Use it to derive facet memberships and answer facts. Anti-hallucination facts, completeness memberships, and conflicting-evidence precedence are later extensions. This is the clinical analogue of PhantomWiki's solver-derived answers \citep{gong2025phantomwiki}.

\stage{Stage 4}{Render Fact Rows into Text Chunks via LLM}

Example rendered chunk:

\begin{lstlisting}[style=promptstyle]
The patient, older than 75, was treated for encephalitis with a 14-day
course of acyclovir. At discharge, orientation had improved, but gait
instability persisted. The patient was transferred to inpatient
rehabilitation for mobility and balance training.
\end{lstlisting}

Store the generated text together with \field{fact\_ids}, \field{facet\_ids}, \field{query\_id}, \field{cluster\_id}, and \field{split}. The output of this stage is \field{chunks.parquet}.

Use template-first, LLM-second rendering: create a deterministic draft from the fact row, then ask the LLM to rewrite it under strict \field{must\_mention}/\field{must\_not\_mention} constraints. Re-extract facts after rendering and reject chunks whose extracted facts disagree with the source row. This addresses PhantomWiki's warning that LLM rephrasing can introduce factual errors \citep{gong2025phantomwiki}.

\stage{Stage 5}{Create Role-Aware Gold Clusters}

Example for \field{Q\_001}:

\begin{itemize}[noitemsep, topsep=0pt]
	\item \field{C1}: encephalitis + uncomplicated diabetes + treatment duration, 18 chunks, \field{cluster\_role=dominant\_gold};
	\item \field{C2}: encephalitis + uncomplicated diabetes + rehabilitation outcome, 8 chunks, \field{cluster\_role=complementary\_gold};
	\item \field{C3}: encephalitis + age $>75$ + treatment duration, 8 chunks, \field{cluster\_role=complementary\_gold};
	\item \field{C4}: encephalitis + age $>75$ + rehabilitation outcome, 8 chunks, \field{cluster\_role=complementary\_gold}.
\end{itemize}

The dominant cluster is larger or more lexically similar to the surface query; this creates the top-$k$ failure mode. High similarity inside \field{C1} is allowed, high similarity between \field{C1} and \field{C2} is suspicious, and a chunk that supports several required facets at once should usually be rejected.

Optionally create \field{completeness\_cluster} groups of 4--10 chunks where every chunk contributes a necessary subclaim and no single chunk contains the full answer. This adapts the completeness-cluster construction in EnterpriseRAG-Bench \citep{sun2026enterpriseragbench}.

\stage{Stage 6}{Generate Hard Distractor Chunks}

For each query plan, generate non-gold chunks that are plausible and semantically close, but do not answer the required facets. These are needed to distinguish coverage from generic diversity.

Useful distractor types:

\begin{itemize}[noitemsep, topsep=0pt]
	\item same condition, wrong subgroup: encephalitis without diabetes and not age $>75$;
	\item same subgroup, wrong condition: elderly stroke patients with rehabilitation outcomes;
	\item same treatment vocabulary, wrong disease: acyclovir duration for shingles;
	\item same outcome vocabulary, wrong clinical axis: rehabilitation outcome after pneumonia;
	\item dispersion trap: semantically different clinical text that is novel but does not support any required facet;
	\item broad background: general encephalitis treatment description;
	\item conflicting near-duplicate: an older note and a later note disagree on a duration or outcome.
\end{itemize}

Store these chunks in \field{chunks.parquet} with \field{distractor\_type} set and no gold qrel for the target query. This is the practical adaptation of diversity-agent ideas from RAG dataset generation \citep{driouich2025diverseprivate}. Each distractor has a reason why it is close to the query and a reason why it is not gold. This poses a challenge for MMR.

\stage{Stage 7}{Generate the Natural Query and Gold Answer via LLMs}

Example:

\begin{lstlisting}[style=promptstyle]
For patients diagnosed with encephalitis or myelitis, how do treatment durations
and rehabilitative outcomes differ between patients with uncomplicated diabetes
and patients older than 75?
\end{lstlisting}
Hidden subqueries can be stored for diagnostics, but the main retrieval experiment should use the single natural query \citep{abbasiantaeb2026multiaspect,yue2021cliniqg4qa,lin2025scirgen}.

Use a finite set of query templates before asking the LLM for fluent wording. Templates should cover subgroup comparison, constrained cohort questions, conflicting evidence, completeness, outcome synthesis, and information-not-found cases. This combines EnterpriseRAG-Bench's question-family design \citep{sun2026enterpriseragbench} with PhantomWiki's grammar-controlled generation \citep{gong2025phantomwiki}.

The answer is also produced at this stage, during benchmark construction, before embedding and before any retrieval method is evaluated. The source of truth is the hidden scaffold: \field{query\_plans.parquet}, \field{clinical\_facts.parquet}, and the fact-to-chunk mapping stored in \field{chunks.parquet}.

Use two separate prompts or procedures:

\begin{enumerate}[noitemsep, topsep=0pt]
	\item \textbf{Query prompt.} Given only the hidden plan, write the user-facing \field{query\_text}. Do not expose the answer wording.
	\item \textbf{Answer prompt.} Given the required facets, the answer outline, and the structured fact rows grouped by facet, write a canonical answer. The prompt may receive \field{fact\_ids} and their derived \field{chunk\_ids} for citation, but it should not receive distractor chunks or retrieval rankings.
\end{enumerate}

Example answer object:

\begin{lstlisting}[style=promptstyle]
query_id: Q_001
answer_text: >
  In the uncomplicated-diabetes subgroup, treatment was generally shorter,
  with acyclovir courses around 7--10 days, and rehabilitation needs were
  usually limited to home therapy or near-baseline mobility. In patients
  older than 75, treatment was more often extended to about 14 days, and
  discharge commonly involved inpatient rehabilitation because gait
  instability or persistent weakness remained. Overall, the elderly subgroup
  shows longer treatment and worse functional recovery than the diabetes
  subgroup.
facet_answer_summaries:
  F1: diabetes subgroup treatment duration pattern
  F2: diabetes subgroup rehabilitation outcome pattern
  F3: age_over_75 subgroup treatment duration pattern
  F4: age_over_75 subgroup rehabilitation outcome pattern
answer_claims:
  - claim_id: A1
    facet_id: F1
    support_fact_ids: [FACT_010, FACT_011]
    support_chunk_ids: [CH_020, CH_021]
  - claim_id: A2
    facet_id: F2
    support_fact_ids: [FACT_018, FACT_019]
    support_chunk_ids: [CH_031, CH_032]
\end{lstlisting}

This is stored in \field{gold\_answers.parquet}. \\

For constrained, conflicting, and hard-distractor questions, also store \field{anti\_hallucination\_facts}: negative claims derived from plausible distractors that should not appear in the answer. These are useful for answer validation and follow EnterpriseRAG-Bench's treatment of distractor-derived false facts \citep{sun2026enterpriseragbench}.

\stage{Stage 8}{Create Qrels Directly from the Hidden Mapping}

Create \field{qrels.parquet} from the known mapping between queries, chunks, facts, and facets.

Example qrels:

\begin{lstlisting}[style=promptstyle]
query_id  chunk_id  facet_id  relevance_grade  fact_ids
Q_001     CH_017    F3        1                [FACT_001]
Q_001     CH_017    F4        1                [FACT_001]
Q_001     CH_088    none      0                []
\end{lstlisting}

\stage{Stage 9}{Validate the Generated Objects}

Run simple automatic checks before embedding:

\begin{itemize}[noitemsep, topsep=0pt]
	\item every query has 3--5 required facets;
	\item every required facet has enough gold chunks;
	\item every \field{cluster\_id} has exactly one \field{cluster\_role};
	\item every generated chunk can be traced back to fact IDs;
	\item chunks do not mention forbidden facts;
	\item distractors do not accidentally answer a gold facet;
	\item BM25/vector/optional agentic pooling does not reveal non-gold chunks that should be qrels;
	\item anti-hallucination facts are not supported by gold chunks;
	\item no chunk answers the entire query unless it is part of an easy ablation;
	\item the canonical answer uses every required facet at least once;
\end{itemize}

Multi-retriever pooling is used as a collision detector, not as the main source of gold truth. If pooling finds required non-gold evidence, regenerate, repair qrels only when the hidden schema agrees, or reject the query \citep{sun2026enterpriseragbench}.

\stage{Stage 10}{Embed and Filter by Retrieval Geometry}

Embed all chunks and queries. These are benchmark-construction checks, so they are allowed to use the hidden structure of the generated data: \field{query\_id}, \field{facet\_ids}, \field{cluster\_id}, \field{cluster\_role}, \field{distractor\_type}, and \field{qrels.parquet}. These fields are used to validate the dataset, but they are not exposed to the retrieval methods.

Separate two different questions:

\begin{itemize}[noitemsep, topsep=0pt]
	\item \textbf{Query-local structure, primary generation check.} Use only the chunks generated for the current query. This checks whether the synthetic item itself has the intended facet and cluster structure.
	\item \textbf{Full-corpus visibility, primary retrieval check if full-corpus evaluation is reported.} Use the full \field{chunks.parquet} corpus. This does not prove the local structure; it only checks whether the gold facets remain visible in the first-stage top-$N$ pool when chunks from other query plans also compete.
\end{itemize}

Query-local checks:
\begin{itemize}[noitemsep, topsep=0pt]
	\item each required facet has enough gold chunks;
	\item intra-facet similarity is high enough that each facet forms a usable cluster;
	\item inter-facet similarity is lower than intra-facet similarity, so facets do not collapse into one cluster;
	\item the dominant gold cluster is larger or more query-similar than the complementary gold clusters;
	\item hard distractors are close to the query or to a gold facet, but have no gold qrel for the target query;
	\item local top-10 nearest-neighbor density, in-facet similarity, cross-facet similarity, and distractor proximity match predeclared thresholds.
\end{itemize}

These geometry diagnostics follow EnterpriseRAG-Bench's use of local-density, in-cluster, and cross-cluster similarity to characterize synthetic RAG corpora \citep{sun2026enterpriseragbench}.

Full-corpus visibility checks, if the main evaluation uses full-corpus retrieval:
\begin{itemize}[noitemsep, topsep=0pt]
	\item each required facet has at least one gold chunk in the full-corpus top-$N$ pool;
	\item the top-$N$ pool contains hard distractors or background negatives, so evaluation is not trivial;
	\item chunks from other query plans do not completely swamp the query's gold evidence.
\end{itemize}

Queries that fail these checks can be regenerated or excluded.

\stage{Stage 11}{Evaluate Retrieval Methods}

The retrieval corpus is part of the benchmark definition, so different corpus choices should be reported as different settings.

\begin{itemize}[noitemsep, topsep=0pt]
	\item \textbf{Full-corpus retrieval, primary setting.} uses the full generated corpus with all the data points belonging to all the queries to realistically test the RAG system. The corpus contains the gold chunks for the current query, its hard distractors, and all chunks generated for the other query plans in the same split.
	\item \textbf{Query-local retrieval, diagnostic ablation.} Run the same methods on only the chunks generated for the current query: gold chunks, hard distractors, and query-specific background negatives. This tests reranking behavior after the query-local universe has already been isolated.
	\item \textbf{Lexical and hybrid sanity baselines.} Optionally include BM25 and hybrid retrieval because EnterpriseRAG-Bench shows that lexical retrieval can outperform vector retrieval when domain-specific terms and identifiers matter \citep{sun2026enterpriseragbench}.
\end{itemize}

\section{Data splits}

\begin{itemize}[]
	\item \textbf{Validation split}: used after the generation pipeline is fixed. We use it to choose retrieval hyperparameters, especially MMR and facility-location $\lambda$ values
	\item \textbf{Test split}: final evaluation
\end{itemize}

\section{Prompt Families}

\begingroup
\renewcommand{\arraystretch}{1.2}
\begin{center}
	\begin{tabular}{@{}p{2.2cm}p{4cm}p{8.2cm}@{}}
		\toprule
		\textbf{Stage} & \textbf{Prompt family} & \textbf{Purpose} \\
		\midrule
		Stage 1 & Optional catalogue refinement & Usually no LLM is needed. If used, ask for plausible conditions, subgroups, treatments, outcomes, and forbidden combinations from a small seed catalogue. The output should be reviewed manually before use. \\
		Stage 2 & Query-plan prompt & Given a catalogue row, produce a hidden plan: primary condition, subgroups, axes, required facets, template/logical form, cluster roles, and difficulty fields. This prompt does not write the final user query. \\
		Stage 3 & Structured fact-row prompt & Given one required facet and its cluster role, generate structured synthetic fact rows with \field{must\_mention}, \field{must\_not\_mention}, metadata fields, and facet IDs. \\
		Stage 4 & Chunk-rendering prompt & Convert one or more fact rows into short clinical-style evidence chunks while preserving the assigned facts and avoiding forbidden facts. \\
		Stage 6 & Distractor-generation prompt & Given a query plan and distractor type, generate plausible non-gold chunks that are close to the query but fail a specific required facet. Store the explicit \field{distractor\_reason}. \\
		Stage 7 & Query and answer prompts & Generate the natural user query from the hidden template/logical form, then generate the canonical answer from grouped structured facts and fact-to-chunk provenance. \\
		Stage 9 & Validation prompts & Use LLM-as-a-judge prompts and multi-retriever pooling to check schema compliance, forbidden-fact leakage, distractor leakage, clinical plausibility, accidental qrel collisions, and whether the canonical answer uses all required facets. These prompts can reject items. \\
		\bottomrule
	\end{tabular}
\end{center}
\endgroup

\section{Metrics}

\begingroup
\renewcommand{\arraystretch}{1.25}
\begin{center}
	\begin{tabular}{@{}p{4cm}p{10.4cm}@{}}
		\toprule
		\textbf{Metric} & \textbf{Purpose} \\
		\midrule
		Facet Coverage@k & Fraction of required facets represented by at least one retrieved gold chunk. In this reduced pipeline each required facet is implemented as one gold cluster, so this is also cluster coverage. Primary metric. \\
		Weighted Facet Coverage@k & Average fraction of pool-visible gold chunks retrieved per facet. \\
		Gold Precision (GP) & Fraction of retrieved chunks that are gold for any facet. \\
		Gold Recall (GR) & Fraction of pool-visible gold chunks retrieved. \\
		Dominant-Cluster Concentration@k & Fraction of selected chunks from the largest or most query-similar gold cluster. Shows top-$k$ redundancy. \\
		Distractor Rate@k & Fraction of selected chunks that are hard distractors. Shows MMR failure modes. \\
		Answer Fact Completeness & Fraction of canonical answer facts whose supporting facet has been retrieved. \\
		Invalid Extra Chunks & Count of selected chunks that are neither gold nor explicitly valid background. \\
		Jaccard vs. Top-$k$ & Whether reranking materially differs from the baseline. \\
		\bottomrule
	\end{tabular}
\end{center}
\endgroup

\section{Benchmark Objects}

\begingroup
\renewcommand{\arraystretch}{1.25}
\begin{center}
	\begin{tabular}{@{}p{4.4cm}p{10.2cm}@{}}
		\toprule
		\textbf{Artifact} & \textbf{Purpose} \\
		\midrule
		\field{clinical\_facts.parquet} & hidden structured source of truth: condition, demographics, comorbidities, interventions, durations, outcomes, complications, and facet IDs. \\
		\field{chunks.parquet} & retrieval corpus: \field{chunk\_id}, \field{fact\_ids}, \field{facet\_ids}, \field{cluster\_id}, \field{distractor\_type}, and text. \\
		\field{queries.parquet} & query text, condition, query type, difficulty, split, and number of facets. \\
		\field{gold\_answers.parquet} & canonical answer object: \field{answer\_text}, facet-level answer summaries, answer claims, supporting fact IDs, and supporting chunk IDs. \\
		\field{query\_plans.parquet} & hidden plan for each query: selected facets, answer outline, split, difficulty, template/logical form, and generation constraints. \\
		\field{qrels.parquet} & Gold evidence: query--chunk--facet labels, relevance grade, support type, fact IDs, and optional span offsets. \\
		\bottomrule
	\end{tabular}
\end{center}
\endgroup

\newpage
\begingroup
\footnotesize
\begin{multicols}{2}
	\bibliographystyle{plainnat-boldtitles}\bibliography{bibliography}

\end{multicols}
\endgroup

\end{document}
